%============================================================
% PREÁMBULO OFICIAL FLUIDODINAMICA_BASE
%============================================================

\documentclass[12pt, a4paper]{article}

%------------------------------------------------------------
% Idioma, codificación y tipografía
%------------------------------------------------------------
\usepackage[spanish]{babel}
\usepackage[utf8]{inputenc}
\usepackage[T1]{fontenc}
\usepackage{lmodern}        % Tipografía moderna, clara y apta para PDF
\usepackage{microtype}      % Optimiza microtipografía (ideal para libros y apuntes)

%------------------------------------------------------------
% Matemática
%------------------------------------------------------------
\usepackage{amsmath, amssymb, amsfonts}
\usepackage{bm}             % Vectores y tensores en negrita: \bm{u}
\usepackage{mathtools}      % Extensiones de amsmath
\usepackage{siunitx}        % Notación correcta de unidades SI
\sisetup{
  locale = DE,            % usa coma decimal si querés; cambiar a US para punto
  per-mode = symbol,
  exponent-product = \cdot,
  separate-uncertainty = true
}

%------------------------------------------------------------
% Gráficos y figuras
%------------------------------------------------------------
\usepackage{graphicx}
\usepackage{float}
\usepackage{subcaption}     % Figuras con subfiguras
\usepackage{wrapfig}

%------------------------------------------------------------
% Diseño de página
%------------------------------------------------------------
\usepackage{geometry}
\geometry{
  a4paper,
  left=25mm,
  right=25mm,
  top=25mm,
  bottom=25mm
}

\usepackage{setspace}       % Control de espaciado (opcional)
\onehalfspacing             % Espaciado 1.5 (queda excelente para apuntes)

%------------------------------------------------------------
% Tablas profesionales
%------------------------------------------------------------
\usepackage{booktabs}       % Líneas horizontales profesionales
\usepackage{multirow}
\usepackage{tabularx}
\usepackage{longtable}

%------------------------------------------------------------
% Enlaces
%------------------------------------------------------------
\usepackage[hidelinks]{hyperref}  % Enlaces sin recuadros

%------------------------------------------------------------
% Colores y cajas
%------------------------------------------------------------
\usepackage{xcolor}
\usepackage{tcolorbox}
\tcbset{
  colback = white,
  colframe = black,
  arc = 2mm,
  boxrule = 0.6pt
}

%------------------------------------------------------------
% Ambientes personalizados
%------------------------------------------------------------

% Entorno para ejercicios
\newenvironment{ejercicio}[1][]{
  \begin{tcolorbox}[title={\textbf{Ejercicio #1}}, colback=gray!5]
}{
  \end{tcolorbox}
}

% Entorno para soluciones
\newenvironment{solucion}{
  \begin{tcolorbox}[title={\textbf{Solución}}, colback=blue!3]
}{
  \end{tcolorbox}
}

% Para dejar comentarios o notas al margen
\newcommand{\nota}[1]{\textcolor{blue}{\textbf{#1}}}

%============================================================
% FIN DEL PREÁMBULO
%============================================================
\begin{document}

%============================================================
% Ejercicio c.2 – Medidor de orificio en la descarga de una bomba
%============================================================

\subsection*{Ejercicio c.2: Placa de orificio en la línea de descarga}

La instalación de la figura se utiliza para transportar agua desde una laguna hasta un
tanque intermedio para distribución. En la línea de conducción se ha instalado un
medidor de orificio de borde a escuadra con tomas de esquina. 

Cuando el nivel en el tanque es el indicado en la figura, la lectura del manómetro
diferencial conectado al orificio es de $18\,\text{cm}$. A su vez, un manómetro ubicado
en la descarga de la bomba indica una presión de $2\,\text{bar}$.

Determinar la longitud de la línea de descarga de la bomba.

\medskip

\noindent
\textbf{Datos adicionales:}
\begin{itemize}
  \item Agua: $\rho = 997\,\text{kg/m}^3$, $\mu = 9{,}7\times 10^{-4}\,\text{Pa\,s}$.
  \item Cañería: IPS Sch 40, diámetro nominal $1''$.
        Diámetro interno: $D \simeq 0{,}0266\,\text{m}$.
        Rugosidad absoluta: $\varepsilon = 4{,}5\times 10^{-5}\,\text{m}$.
  \item Línea de succión: 1 válvula de retención, 1 codo $90^\circ$ (no se usan en este ítem).
  \item Línea de descarga: 1 válvula globo, 2 codos $90^\circ$, 1 orificio de borde a escuadra.
  \item Propiedades del fluido manométrico: $\rho_\text{man} = 13500\,\text{kg/m}^3$ (Hg).
  \item Diámetro del orificio: $d_o = 2{,}0\,\text{cm} = 0{,}020\,\text{m}$.
  \item Lectura del manómetro diferencial del orificio: $\Delta h = 18\,\text{cm}$.
  \item La pérdida de presión permanente del orificio es el $40\,\%$ de la
        diferencia de presión medida por el manómetro asociado al orificio.
  \item El tanque intermedio está venteado a la atmósfera.
  \item Diferencia de cota entre el eje de la cañería de descarga y el nivel
        libre del tanque: $z_3 - z_2 = 2{,}0\,\text{m}$.
  \item Manómetro en la descarga de la bomba: $P_2 = 2\,\text{bar}$ (presión manométrica).
\end{itemize}

\noindent
\textbf{Incógnita:} longitud de la línea de descarga $L$ [m].

%------------------------------------------------------------
\subsubsection*{1. Hipótesis}

\begin{enumerate}
  \item Régimen estacionario, agua incompresible, propiedades constantes.
  \item Se desprecia la velocidad en la superficie libre del tanque ($V_3\simeq 0$).
  \item Tanto el manómetro en la descarga de la bomba como el del orificio
        miden presiones manométricas.
  \item El diámetro en la descarga es constante, por lo que la velocidad media $V$
        es la misma a lo largo de toda la línea (incluyendo la sección del orificio,
        salvo en la vena contraída).
  \item El orificio se modela según el teórico de ``medidores por restricción de área'':
        $\;Q = C_D A_o \sqrt{2\Delta P/[\rho(1-\beta^4)]}\;$ para un orificio de borde a escuadra.
  \item La altura geométrica relevante es la diferencia entre el eje de la cañería de descarga
        y el nivel libre del tanque: $z_3 - z_2 = 2\,\text{m}$.
\end{enumerate}

%------------------------------------------------------------
\subsubsection*{2. Modelo del orificio y caudal en la línea}

\paragraph{2.1. Diferencia de presión en el orificio.}

Con tomas de esquina al mismo nivel y agua como fluido de proceso, el salto de presión
entre las secciones 1 (antes del orificio) y 2 (después del orificio) viene dado por
la lectura del manómetro diferencial con fluido manométrico de densidad
$\rho_\text{man}$:

\[
\Delta P_{\text{man}} = P_1 - P_2 = (\rho_\text{man} - \rho)\,g\,\Delta h.
\]

Sustituyendo numéricamente:

\[
\Delta P_{\text{man}}
= (13500 - 997)\,\text{kg/m}^3 \cdot 9{,}81\,\text{m/s}^2 \cdot 0{,}18\,\text{m}
\approx 2{,}21\times 10^4\,\text{Pa}.
\]

\paragraph{2.2. Ecuación de orificio.}

Para el orificio de borde a escuadra con flujo turbulento, tus notas de clase
indican un coeficiente de descarga típico $C_D \simeq 0{,}62$ cuando
$\text{Re}>3\times 10^4$.

La relación de diámetros es

\[
\beta = \frac{d_o}{D} = \frac{0{,}020}{0{,}0266} \approx 0{,}75,
\qquad
1-\beta^4 \approx 0{,}68.
\]

El área del orificio es

\[
A_o = \frac{\pi d_o^2}{4}
    = \frac{\pi (0{,}020)^2}{4}
    = 3{,}14\times 10^{-4}\,\text{m}^2.
\]

La ecuación reducida de orificio (análogo del Venturi) es

\[
Q = C_D\,A_o\,
    \sqrt{\frac{2\,\Delta P_{\text{man}}}{\rho (1-\beta^4)}}.
\]

Sustituyendo:

\[
Q \approx 0{,}62 \times 3{,}14\times 10^{-4}
    \sqrt{\frac{2\cdot 2{,}21\times 10^4}{997\cdot 0{,}68}}
   \;\approx\; 1{,}6\times 10^{-3}\,\text{m}^3/\text{s}.
\]

\[
\boxed{Q \approx 1{,}6\times 10^{-3}\,\text{m}^3/\text{s}
      \;\; (\approx 5{,}7\,\text{m}^3/\text{h}).}
\]

La velocidad media en la cañería de descarga (diámetro $D$) es

\[
A = \frac{\pi D^2}{4}
  = \frac{\pi (0{,}0266)^2}{4}
  \approx 5{,}56\times 10^{-4}\,\text{m}^2,
\]

\[
V = \frac{Q}{A}
  \approx \frac{1{,}6\times 10^{-3}}{5{,}56\times 10^{-4}}
  \approx 2{,}8\,\text{m/s}.
\]

El número de Reynolds en la cañería es

\[
\text{Re} = \frac{\rho V D}{\mu}
          \approx \frac{997 \cdot 2{,}8 \cdot 0{,}0266}{9{,}7\times 10^{-4}}
          \sim 7\times 10^4,
\]

lo que valida el uso de $C_D\simeq 0{,}62$.

%------------------------------------------------------------
\subsubsection*{3. Factor de fricción en la descarga}

Usando la correlación de Churchill/Colebrook con $\varepsilon = 4{,}5\times 10^{-5}\,\text{m}$,
$\text{Re}\sim 7\times 10^4$ y $D=0{,}0266\,\text{m}$, se obtiene un factor de fricción
de Darcy del orden

\[
c_f \simeq 0{,}024\text{--}0{,}025.
\]

En lo que sigue se tomará un valor típico
\[
c_f \approx 0{,}024,
\]
acorde con tus tablas y con el cálculo por Churchill en la planilla correspondiente.

La energía cinética específica en la línea es

\[
\frac{V^2}{2g}
= \frac{(2{,}8)^2}{2\cdot 9{,}81}
\approx 0{,}41\,\text{m}.
\]

%------------------------------------------------------------
\subsubsection*{4. Balance de energía entre la descarga de la bomba y el tanque}

Consideramos la sección $2$ en la descarga de la bomba (donde está el manómetro) 
y la sección $3$ en el nivel libre del tanque. Tomamos el eje de la cañería de descarga
como referencia $z_2 = 0$, de modo que $z_3 = 2{,}0\,\text{m}$.

El balance de energía entre $2$ y $3$ (sin máquinas entre ambas secciones) es

\[
\frac{P_2}{\rho g} + z_2 + \frac{V^2}{2g}
=
\frac{P_3}{\rho g} + z_3 + \frac{V_3^2}{2g} + h_{L,23},
\]

donde $h_{L,23}$ son las pérdidas totales entre 2 y 3. Como el tanque está venteado,
$P_3 = P_{\text{atm}}$ y el manómetro en 2 mide presión manométrica, se tiene
$P_2/(\rho g)$ directamente como carga de presión en 2:

\[
\frac{P_2}{\rho g}
= \frac{2\times 10^5}{997\cdot 9{,}81}
\approx 20{,}4\,\text{m}.
\]

Además, $V_3 \simeq 0$, por lo que

\[
h_{L,23}
= \frac{P_2}{\rho g} + 0 + \frac{V^2}{2g}
- \left(0 + 2{,}0 + 0\right)
\approx 20{,}4 + 0{,}41 - 2{,}0
\approx 18{,}9\,\text{m}.
\]

Es decir, la bomba entrega una carga tal que, entre su descarga y el tanque,
se pierden aproximadamente $18{,}9\,\text{m}$ de columna de agua.

%------------------------------------------------------------
\subsubsection*{5. Descomposición de las pérdidas $h_{L,23}$}

Descomponemos las pérdidas en:

\[
h_{L,23} = h_{f,L} + h_{\text{loc,desc}} + h_{\text{perm, orif}},
\]

donde:

\begin{itemize}
  \item $h_{f,L}$: pérdida por fricción distribuida en la longitud recta $L$,
        \[
        h_{f,L} = c_f\,\frac{L}{D}\,\frac{V^2}{2g}.
        \]
  \item $h_{\text{loc,desc}}$: pérdidas localizadas en la descarga \emph{excepto}
        la pérdida permanente del orificio (válvula globo, codos $90^\circ$ y
        descarga sumergida).
  \item $h_{\text{perm, orif}}$: pérdida permanente asociada al orificio
        (dato del enunciado).
\end{itemize}

\paragraph{5.1. Pérdida permanente del orificio.}

La pérdida permanente es el $40\,\%$ de la diferencia de presión medida:

\[
h_{\text{perm, orif}}
= 0{,}40\,\frac{\Delta P_{\text{man}}}{\rho g}
\approx 0{,}40\,
\frac{2{,}21\times 10^4}{997\cdot 9{,}81}
\approx 0{,}90\,\text{m}.
\]

\paragraph{5.2. Pérdidas localizadas en válvula, codos y descarga sumergida.}

De acuerdo con las tablas de Crane (usadas en la asignatura), expresadas como
longitudes equivalentes:

\[
\left(\frac{L_e}{D}\right)_\text{globo} \approx 340,\qquad
\left(\frac{L_e}{D}\right)_\text{codo} \approx 30,
\]

de modo que para 1 válvula globo y 2 codos:

\[
\left(\frac{L_e}{D}\right)_\text{total desc}
= 340 + 2\cdot 30 = 400.
\]

La pérdida equivalente de estos accesorios es

\[
h_{\text{valv+codos}}
= c_f\left(\frac{L_e}{D}\right)_\text{total desc}\frac{V^2}{2g}
= c_f\,400\,\frac{V^2}{2g}.
\]

Además, la descarga sumergida produce una pérdida adicional asociada a la
pérdida completa de la energía cinética:

\[
h_{\text{descarga sumergida}} = K\,\frac{V^2}{2g},
\qquad K = 1.
\]

Por tanto,

\[
h_{\text{loc,desc}}
= c_f\,400\,\frac{V^2}{2g} + 1\cdot \frac{V^2}{2g}.
\]

\paragraph{5.3. Cálculo numérico.}

Con $c_f \approx 0{,}024$ y $V^2/(2g)\approx 0{,}41\,\text{m}$:

\[
h_{\text{valv+codos}}
\approx 0{,}024\cdot 400\cdot 0{,}41
\approx 3{,}9\,\text{m},
\]

\[
h_{\text{descarga sumergida}} \approx 1\cdot 0{,}41 \approx 0{,}4\,\text{m},
\]

\[
h_{\text{loc,desc}} \approx 3{,}9 + 0{,}4 \approx 4{,}3\text{--}4{,}5\,\text{m}.
\]

\paragraph{5.4. Pérdida por fricción en la tubería recta y longitud $L$.}

Del balance global:

\[
h_{f,L}
= h_{L,23} - h_{\text{loc,desc}} - h_{\text{perm, orif}}
\approx 18{,}9 - 4{,}4 - 0{,}9
\approx 13{,}6\,\text{m}.
\]

Pero

\[
h_{f,L} = c_f\,\frac{L}{D}\,\frac{V^2}{2g}
\quad\Rightarrow\quad
L = \frac{h_{f,L}\,D}{c_f}\,\frac{2g}{V^2}.
\]

Sustituyendo $h_{f,L}\approx 13{,}6\,\text{m}$,
$c_f \approx 0{,}024$, $D = 0{,}0266\,\text{m}$,
$V^2/(2g)\approx 0{,}41\,\text{m}$:

\[
L \approx
\frac{13{,}6\cdot 0{,}0266}{0{,}024}\cdot
\frac{1}{0{,}41}
\;\approx\; 3{,}8\times 10^{1}\,\text{m}.
\]

\[
\boxed{L \approx 3{,}8\times 10^{1}\,\text{m}
      \;\; (\text{longitud de descarga del orden de } 38\,\text{m}).}
\]

%------------------------------------------------------------
\subsubsection*{6. Comentario final}

\begin{itemize}
  \item El orificio se ha modelado exactamente como en el teórico de medidores
        por restricción de área, con $C_D\simeq 0{,}62$ y el factor geométrico
        $(1-\beta^4)$.
  \item La diferencia de presión en el orificio se obtuvo a partir del manómetro
        diferencial con fluido de alta densidad (Hg), y se utilizó la fracción
        $40\,\%$ para la pérdida permanente, como indica el enunciado.
  \item En el balance de energía 2--3 se incluyeron:
        \begin{enumerate}
          \item la elevación geométrica $(z_3-z_2)$,
          \item la pérdida distribuida en la longitud $L$,
          \item las pérdidas localizadas de la válvula globo y los dos codos,
          \item y, muy importante, la pérdida con $K=1$ asociada a la 
                \emph{descarga sumergida}, que representa la disipación de
                toda la energía cinética del chorro al ingresar al tanque.
        \end{enumerate}
  \item El resultado $L \approx 38\,\text{m}$ es coherente con la orden de 
        magnitud de las pérdidas calculadas y con el régimen turbulento de 
        escoamiento en una cañería de $1''$.
\end{itemize}

\end{document}